\documentclass[12pt,twoside]{article}
\usepackage[utf8]{inputenc}
\usepackage[a4paper,width=160mm,top=25mm,bottom=25mm,bindingoffset=6mm]{geometry}
\usepackage{tikz}
\usetikzlibrary{calc}
\usepackage{eso-pic}
\usepackage{setspace}
\usepackage{lipsum}
\usepackage[english]{babel}
\usepackage{array}
\usepackage{tabu}
\usepackage{multirow}
\usepackage{gensymb}
\usepackage{rotating}
\usepackage{caption}
\usepackage{subcaption}
\usepackage{hanging}
\usepackage{changepage}
\usepackage{subcaption}
\usepackage{lscape}
\usepackage{chngcntr}
\usepackage{float}
\floatstyle{plaintop}
\restylefloat{table}
\usepackage{caption}
\usepackage{sectsty}
\usepackage{enumitem}
\usepackage{hyperref}
\hypersetup{
    colorlinks,
    linkcolor={blue},
    citecolor={blue},
    urlcolor={blue}
}
\usepackage{amsmath}
\usepackage{mathtools}
\usepackage{amssymb}
\newcommand{\jump}[1]{\ensuremath{[\![#1]\!]} }

\title{{\textbf{Discontinuous Galerkin weak formulation for geometrically exact Kirchhoff-Love beams}}}
\author{}
\date{}

\sectionfont{\fontsize{12}{12}\selectfont}

\begin{document}

\maketitle

\noindent
From the law of conservation of linear and angular momentum, the strong form static equilibrium equations can be written follows:

\begin{gather}
  \boldsymbol{f^{\prime}} + \boldsymbol{\tilde{f}} = \boldsymbol{0}\label{eq1} \\
  \boldsymbol{m^{\prime}} + \boldsymbol{r^{\prime}} \times \boldsymbol{f} + \boldsymbol{\tilde{m}} = \boldsymbol{0}\label{eq2}
\end{gather}

\noindent
Applying the Kirchhoff constraint (negligible shear deformation) and considering isotropic beams in a torsion-free situation, equations~\ref{eq1} and~\ref{eq2} can be reduced to the following system of three differential equations:

\begin{gather}
  \boldsymbol{f^{\prime}_{||}} + \Bigg[\frac{\boldsymbol{r^{\prime}}}{||\boldsymbol{r^{\prime}}||^{2}} \times \big(\boldsymbol{m^{\prime}} + \boldsymbol{\tilde{m}_{\perp}}\big)\Bigg]^{\prime} + \boldsymbol{\tilde{f}} = \boldsymbol{0}\label{eq3}
\end{gather}

\noindent
Multiplying equations~\ref{eq3} by $\delta\boldsymbol{r}$ \& performing integrating by parts along with Gauss-Divergence theorem over the elements e = 1, 2, ..., E gives:

\begin{gather}
  \sum_{e=1}^{E} \int_{S_{e}} \delta\boldsymbol{r} \cdot \Bigg[\boldsymbol{f_{||}} + \frac{\boldsymbol{r^{\prime}}}{||\boldsymbol{r^{\prime}}||^{2}} \times \big(\boldsymbol{m^{\prime}} + \boldsymbol{\tilde{m}_{\perp}}\big)\Bigg]^{\prime}ds + \sum_{e=1}^{E} \int_{S_{e}} \delta\boldsymbol{r} \cdot \boldsymbol{\tilde{f}}\:ds = 0
\end{gather}

\begin{gather}
  \sum_{e=1}^{E} \Bigg[\delta\boldsymbol{r} \cdot \Bigg[\boldsymbol{f_{||}} + \frac{\boldsymbol{r^{\prime}}}{||\boldsymbol{r^{\prime}}||^{2}} \times \big(\boldsymbol{m^{\prime}} + \boldsymbol{\tilde{m}_{\perp}}\big)\Bigg]\Bigg]_{\partial S_{e}} + \Bigg[\delta\boldsymbol{r} \cdot \Bigg[\boldsymbol{f_{||}} + \frac{\boldsymbol{r^{\prime}}}{||\boldsymbol{r^{\prime}}||^{2}} \times \big(\boldsymbol{m^{\prime}} + \boldsymbol{\tilde{m}_{\perp}}\big)\Bigg]\Bigg]_{\partial_{N} S} \cdots \notag \\
  - \sum_{e=1}^{E} \int_{S_{e}} \delta\boldsymbol{r^{\prime}} \cdot \Bigg[\boldsymbol{f_{||}} + \frac{\boldsymbol{r^{\prime}}}{||\boldsymbol{r^{\prime}}||^{2}} \times \big(\boldsymbol{m^{\prime}} + \boldsymbol{\tilde{m}_{\perp}}\big)\Bigg]ds + \sum_{e=1}^{E} \int_{S_{e}} \delta\boldsymbol{r} \cdot \boldsymbol{\tilde{f}}\:ds = 0
\end{gather}

\begin{gather}
  \Big[\delta\boldsymbol{r} \cdot \boldsymbol{\overline{f}}\Big]_{\partial_{N} S} - \Big[\jump{\delta\boldsymbol{r} \cdot \boldsymbol{f}}\Big]_{\partial_{I} S} - \int_{S} \delta\boldsymbol{r^{\prime}}\cdot\boldsymbol{f_{||}}\:ds + \int_{S} \delta \boldsymbol{\theta_{\perp}}\cdot\boldsymbol{\tilde{m}_{\perp}}\:ds \cdots \notag \\
  + \sum_{e=1}^{E} \int_{S_{e}} \delta \boldsymbol{\theta_{\perp}}\cdot\boldsymbol{m^{\prime}}\:ds + \int_{S} \delta\boldsymbol{r} \cdot \boldsymbol{\tilde{f}}\:ds = 0 \label{eq6}
\end{gather}

\noindent
where $S = \bigcup\limits_{e=1}^{E}\overline{S}_{e}$ is the discretization of the computational domain, $\overline{S}_{e} = S_{e}\cup\partial S_{e}$ contains the open domain $S_{e}$ and its boundary $\partial S_{e}$, $\partial_{I} S$ contains the inter-element boundaries of the discretization, and $\partial_{N} S$ is the Neumann boundary of the discretization.

\noindent
Performing integration by parts along with Gauss-Divergence theorem again for the term in equation~\ref{eq6} with derivative of moments gives:

\begin{gather}
  \Big[\delta\boldsymbol{r} \cdot \boldsymbol{\overline{f}}\Big]_{\partial_{N} S} - \Big[\jump{\delta\boldsymbol{r} \cdot \boldsymbol{f}}\Big]_{\partial_{I} S} - \int_{S} \delta\boldsymbol{r^{\prime}}\cdot\boldsymbol{f_{||}}\:ds + \int_{S} \delta \boldsymbol{\theta_{\perp}}\cdot\boldsymbol{\tilde{m}_{\perp}}\:ds + \sum_{e=1}^{E} \Big[\delta \boldsymbol{\theta_{\perp}}\cdot\boldsymbol{m}\Big]_{\partial_{I} S} \cdots \notag \\
  + \Big[\delta \boldsymbol{\theta_{\perp}}\cdot\boldsymbol{\overline{m}}\Big]_{\partial_{N} S} - \sum_{e=1}^{E} \int_{S_{e}} \delta \boldsymbol{\theta_{\perp}^{\prime}}\cdot\boldsymbol{m}\:ds + \int_{S} \delta\boldsymbol{r} \cdot \boldsymbol{\tilde{f}}\:ds = 0
\end{gather}

\begin{gather}
  \Big[\delta\boldsymbol{r} \cdot \boldsymbol{\overline{f}}\Big]_{\partial_{N} S} + \Big[\delta \boldsymbol{\theta_{\perp}}\cdot\boldsymbol{\overline{m}}\Big]_{\partial_{N} S}  + \int_{S} \delta\boldsymbol{r} \cdot \boldsymbol{\tilde{f}}\:ds + \int_{S} \delta \boldsymbol{\theta_{\perp}}\cdot\boldsymbol{\tilde{m}_{\perp}}\:ds = \int_{S} \delta\boldsymbol{r^{\prime}}\cdot\boldsymbol{f_{||}}\:ds \cdots \notag \\
  + \int_{S} \delta \boldsymbol{\theta_{\perp}^{\prime}}\cdot\boldsymbol{m}\:ds + \Big[\jump{\delta\boldsymbol{r} \cdot \boldsymbol{f}}\Big]_{\partial_{I} S} + \Big[\jump{\delta \boldsymbol{\theta_{\perp}}\cdot\boldsymbol{m}}\Big]_{\partial_{I} S}
\end{gather}

\begin{gather}
  \Big[\delta\boldsymbol{r} \cdot \boldsymbol{\overline{f}}\Big]_{\partial_{N} S} + \Big[\delta \boldsymbol{\theta_{\perp}}\cdot\boldsymbol{\overline{m}}\Big]_{\partial_{N} S}  + \int_{S} \delta\boldsymbol{r} \cdot \boldsymbol{\tilde{f}}\:ds + \int_{S} \delta \boldsymbol{\theta_{\perp}}\cdot\boldsymbol{\tilde{m}_{\perp}}\:ds = \int_{S} \delta\boldsymbol{r^{\prime}}\cdot\boldsymbol{f_{||}}\:ds \cdots \notag \\
  + \int_{S} \delta \boldsymbol{\theta_{\perp}^{\prime}}\cdot\boldsymbol{m}\:ds + \Big[\jump{\delta\boldsymbol{r}} \cdot \langle\boldsymbol{f}\rangle\Big]_{\partial_{I} S} + \Big[\langle\delta\boldsymbol{r}\rangle \cdot \jump{\boldsymbol{f}}\Big]_{\partial_{I} S} + \Big[\jump{\delta \boldsymbol{\theta_{\perp}}}\cdot\langle\boldsymbol{m}\rangle\Big]_{\partial_{I} S} \cdots \notag \\
  + \Big[\langle\delta \boldsymbol{\theta_{\perp}}\rangle\cdot\jump{\boldsymbol{m}}\Big]_{\partial_{I} S}
\end{gather}

\noindent
Since forces and moments need not be penalized to ensure consistency and adding a quadratic penalty term on positions for stability of the numerical scheme:

\begin{gather}
  \Big[\delta\boldsymbol{r} \cdot \boldsymbol{\overline{f}}\Big]_{\partial_{N} S} + \Big[\delta \boldsymbol{\theta_{\perp}}\cdot\boldsymbol{\overline{m}}\Big]_{\partial_{N} S}  + \int_{S} \delta\boldsymbol{r} \cdot \boldsymbol{\tilde{f}}\:ds + \int_{S} \delta \boldsymbol{\theta_{\perp}}\cdot\boldsymbol{\tilde{m}_{\perp}}\:ds = \int_{S} \delta\boldsymbol{r^{\prime}}\cdot\boldsymbol{f_{||}}\:ds \cdots \notag \\
  + \int_{S} \delta \boldsymbol{\theta_{\perp}^{\prime}}\cdot\boldsymbol{m}\:ds + \Big[\jump{\delta\boldsymbol{r}} \cdot \langle\boldsymbol{f}\rangle\Big]_{\partial_{I} S} + \Big[\jump{\delta \boldsymbol{\theta_{\perp}}}\cdot\langle\boldsymbol{m}\rangle\Big]_{\partial_{I} S} + \beta\Big[\jump{\delta\boldsymbol{r}}\cdot\jump{\boldsymbol{r}}\Big]_{\partial_{I} S}
\end{gather}

\noindent
where $\beta$ is the penalty parameter.

\end{document}